The integration of blockchain with decentralized federated learning (DFL) has gained significant attention as a mechanism to enhance trust, auditability, and coordination in distributed learning systems. Formery et al.~\cite{BCFramework} provide a comprehensive architectural analysis of Blockchain-Based Decentralized Federated Learning (BDFL), decomposing such systems into layered components including blockchain infrastructure, client interconnection, consensus management, and data transmission. Their work highlights decentralization, immutability, and transparency as key properties that enable trustworthy coordination in open and partially trusted environments. By anchoring model commitments and learning events on-chain, BDFL architectures improve traceability and tamper resistance.

However, these frameworks primarily focus on architectural characterization rather than optimization-aware learning design. The learning layer is typically treated independently from the blockchain layer, with limited consideration of communication topology, data heterogeneity, or aggregation structure. Moreover, most BDFL designs assume a flat peer-to-peer network without structured hierarchy. Our work extends this coordination paradigm by tightly integrating blockchain-based auditability with topology-aware decentralized optimization and multi-level aggregation.

A fundamental challenge in decentralized stochastic gradient descent (D-SGD) is convergence under heterogeneous (non-IID) data distributions. The D-Cliques approach~\cite{Dclique} addresses this issue by organizing nodes into densely connected cliques whose aggregated label distributions approximate the global distribution, while maintaining sparse inter-clique connectivity to reduce communication cost. This topology mitigates optimization bias induced by heterogeneous neighborhoods and improves convergence stability.

Le~Bars et al.~\cite{Refined} further refine the theoretical understanding of D-SGD by introducing the notion of \emph{neighborhood heterogeneity}. Their analysis demonstrates that convergence depends not only on global heterogeneity but also on the mismatch between a node’s data distribution and that of its neighbors. By explicitly coupling topology and heterogeneity, they show that appropriate topology learning can significantly reduce convergence degradation. These works provide strong theoretical and practical motivation for topology-aware decentralized training.

Nevertheless, both D-Cliques and refined topology learning concentrate on improving convergence within a single-level network and do not address system-level concerns such as verifiable coordination, privacy-preserving aggregation across administrative domains, or hierarchical scalability. In contrast, our framework extends topology-aware principles into a structured hierarchy composed of clusters, states, and a nation-level aggregation layer. At the lowest level, nodes are grouped into statistically balanced clusters inspired by D-Cliques. Multiple clusters form a state, and multiple states contribute to nation-level aggregation. This hierarchy enables scalable coordination across organizational boundaries while preserving topology-aware convergence benefits.

Our topology also allowed us to address privacy, which remains critical in decentralized learning. Bonawitz et al.~\cite{SecureAggregation} propose a practical secure aggregation protocol that allows an aggregator to compute only the sum of client updates using pairwise masking and secret sharing, ensuring robustness to dropouts while preventing exposure of individual updates. Although highly effective, this protocol is typically applied in centralized federated learning with a trusted server.

Our framework adapts secure aggregation to the cluster level of a decentralized hierarchy. Within each cluster, a designated cluster head coordinates secure aggregation without accessing individual model updates. At higher levels, aggregation operates on already aggregated cluster models, reducing privacy exposure while enabling scalable learning. This layered aggregation strategy preserves confidentiality while supporting large-scale deployment.

Scalability challenges in blockchain-based federated learning have also motivated off-chain storage solutions. Zhao et al.~\cite{BCIOT} propose leveraging the InterPlanetary File System (IPFS) to store model parameters off-chain, while anchoring content identifiers (CIDs) and hashes on the blockchain. This separation between verification and storage improves scalability and cost efficiency while maintaining integrity guarantees.

However, existing blockchain-FL systems often adopt flat orchestration models and lack structured multi-level aggregation. Our framework incorporates the blockchain/IPFS separation paradigm within a three-level hierarchy: cluster models are securely aggregated and stored in IPFS with blockchain anchoring, state-level merging combines multiple clusters, and nation-level aggregation integrates state models. This yields end-to-end traceability and tamper-resistant model lineage across all levels.

Prior work typically advances one dimension at a time—blockchain coordination~\cite{BCFramework}, topology-aware convergence~\cite{Dclique,Refined}, secure aggregation~\cite{SecureAggregation}, or off-chain storage~\cite{BCIOT}—and leaves the others untouched. Our contribution is to unify these components in a cross-silo setting with stable identities and administrative autonomy. The cluster–state–nation architecture addresses convergence, privacy, auditability, and scalability simultaneously by coupling topology-aware decentralized learning with secure aggregation and blockchain-anchored coordination across all layers, moving blockchain-based decentralized FL closer to practical large-scale deployment.
